\documentclass[main.tex]{subfiles}
\begin{document}
\section{Motive und Prinzipien}
Die Motivation ist die Vermeidung unerwünschter Effekte und die Förderung gewünschter Effekte, wie:
\begin{itemize}
    \item Flexibilität bei sich ändernden Anforderungen
    \item Höhere Kundenzufriedenheit durch kontinuierliches Feedback
    \item Verbesserte Teamarbeit und Kommunikation
    \item Schnellere Lieferzeiten und bessere Qualität: Vermeidung von Zeitverzug 
    \item Vermeidung von Über-Analyse
    \item Vermeidung von Bürokratisierung des Entwicklungsprozesses
\end{itemize}
Die Motivation für agile Methoden ist, die Effizienz und Effektivität der Softwareentwicklung zu steigern, indem man:
\begin{itemize}
    \item Kundenengagement fördert
    \item Risiken frühzeitig identifiziert und mindert
    \item Die Anpassungsfähigkeit an sich ändernde Marktbedingungen verbessert
\end{itemize}
Das Geschieht durch eine Verschlankung der Prozesse: 
\begin{itemize}
    \item Entwurfs- und Planungsphase minimieren
    \item Entwicklung fokussiert auf die Softwarefunktionen, die direkten Wert für den späteren Benutzer haben
    \item Implementierungsphase möglichst frühzeitig starten
    \item Inkrementelle Entwicklung und Bereitstellung
    \item Schnell ausführbare Software für Kundenabstimmung
    \item Frühzeitig und flexibel auf Kundenwünsche eingehen
\end{itemize}
\newpage
\section{Scrum als agiles Prozessmodell}
\subsection*{Es existiert eine große Anzahl agiler Methoden und Vorschläge für agile Prozesse! Scrum ist nur eine davon. }
\begin{myTheorem}{Scrum}{thm:scrum}
Scrum ist ein Rahmenwerk für die agile Softwareentwicklung, basierend auf der hyperproduktiven Produktentwicklung. Es wird engesetzt, wenn Anforderungen sich schnell ändern oder unklar sind und nutzt einen empirischen Ansatz basierend auf der Prozess-Steuerungstheorie. 
\end{myTheorem}
\subsection{Überblick}
\begin{multicols}{2}
\subsubsection*{Die Vorteile}
\begin{itemize}
    \item Der \textbf{Fortschritt} wird nicht durch instabile Anforderungen aufgehalten
    \item Das Produkt wird \textbf{heruntergebrochen} auf eine Menge von verständlichen Einheiten, zu denen die Stakeholder einen Bezug finden können
    \item Kunden erhalten Inkremente \textbf{pünktlich} geliefert und bekommen Feedback darüber, wie das Produkt arbeitet
    \item Die \textbf{Kommunikation} im Team wird verbessert, weil jeder Einsicht in alles hat
    \item Es wird \textbf{Vertrauen} zwischen Kunden und Entwicklern aufgebaut und eine positive Kultur erzeugt
\end{itemize}
\columnbreak
\subsubsection*{Die Nachteile}
\begin{itemize}
    \item Hoher \textbf{Kommunikationsaufwand}
    \item \textbf{Entwurfsphase fehlt} $\Rightarrow$ Architekturentwicklung muss explizit berücksichtigt werden
    \item Eher geeignet für \textbf{kleine Teams}
    \item Herausforderungen bei der \textbf{Vertragsgestaltung}
\end{itemize}
\end{multicols}

\subsubsection*{Empirische Prozessteuerung}
\begin{itemize}
    \item Säulen: Transparenz, Inspektion und Anpassung von Ergebnissen und Aktivitäten
    \item Kombination von iterativen und inkrementellen Vorgehen
    \item Grundidee: Entscheidungsspielraum sinkt von Iteration zu Iteration
    \item Anzahl der messbaren Teilergebnisse steigt stetig
    \item Unschärfe des Projektziels nimmt im Projektverlauf ab
    \item Vorteile:
    \begin{itemize}
        \item Frühe Rückkopplung der Erfahrungen
        \item Reduzierung von Risiken in der Umsetzung
    \end{itemize}
\end{itemize}
\subsubsection*{Grundprinzipien}
\begin{itemize}
    \item Selbstorganisation
    \item Pull-Prinzip
    \item Timebox
    \item Nutzbare Funktionalität
    \item Rahmenwerk
\end{itemize}
\subsubsection*{Rollen}
\begin{itemize}
    \item \textbf{\gls{po}} (Auftraggeber*in) 
    \begin{itemize}
        \item 
    \end{itemize}
    \item \textbf{\gls{sm}}
    \item \textbf{Entwicklungsteam}
\end{itemize}
\subsubsection*{Artefakte}
\begin{myTheorem}{Product Backlog Items}{thm:pbis}
\glspl{pbi} beziehen sich auf alle Anforderungen, die im Product Backlog eines Projekts festgehalten werden. Sie können verschiedene Formen annehmen, darunter:
\begin{itemize}
    \item User Stories
    \item Funktionale und nicht funktionale Anforderungen
    \item Technische Aufgaben
    \item Fehlerbehebungen
\end{itemize}
\end{myTheorem}
\begin{itemize}
    \item \textbf{Product Backlog}
    \begin{itemize}
        \item Anforderungsliste mit Aufgaben
        \item Elemente sind sog. \glspl{pbi}
        \item sollte immer so priorisiert werden, dass die zuerst zu implementierenden Punkte ganz oben auf der Liste stehen
    \end{itemize}
     \item \textbf{Produkt-Inkrement}
    \item \textbf{\acrlong{pbi}}
    \item \textbf{Vision}
    \item \textbf{\gls{spl}}
    \item \textbf{Sprint Backlog}
    \item \textbf{Burndown-Chart}
    \item \textbf{Impediment Backlog}
    \item \textbf{Sprint Goal }
\end{itemize}
\subsubsection*{Ereignisse}
\begin{itemize}
    \item Sprint
    \item Sprint Planning
    \item Daily Scrum
    \begin{itemize}
        \item Kurz
    \end{itemize}
    \item Sprint Review
    \item Sprint Retrospective
\end{itemize}

%\subsubsection*{}

\end{document}